\subsection{Operator Functions}


% **Role:** structure revealed by invariant subspaces and eigenbehavior.
%
% Finite spectral calculus, positive operators, spectral commutativity, and diagonalizable exponentials.

This subsection develops finite spectral functions of normal operators, positive square roots, commutativity through spectral projections, and matrix exponentials for diagonalizable matrices.

\begin{definition} \label{def:normal-spectral-resolution-notation} \hlt{Spectral Resolution Notation}\\
Let $T\in\LL(\mathcal{V})$ be normal on a finite-dimensional complex inner-product space.\\
Define $T=\sum\limits_{\lambda\in\Omega_T}\lambda P_{\lambda}$, where $\Omega_T=\operatorname{spec}(T)$.
\begin{enumerate}[label=(\roman*)]
\setlength{\itemsep}{0pt}
\item \hlt{Spectral Set}: $\Omega_T=\operatorname{spec}(T)$.
\item \hlt{Spectral Projection}: $P_{\lambda}$ is the orthogonal projection onto $E_{\lambda}(T)$.
\item \hlt{Spectral Resolution}: $\{P_{\lambda}\}_{\lambda\in\Omega_T}$, equivalently $T=\sum\limits_{\lambda\in\Omega_T}\lambda P_{\lambda}$.
\end{enumerate}
\end{definition}

\begin{theorem} \label{thm:finite-spectral-functional-calculus-normal} \hlt{Finite Spectral Functional Calculus for Normal Operators}\\
Let $T\in\LL(\mathcal{V})$ be normal, with spectral resolution $T=\sum\limits_{\lambda\in\Omega_T}\lambda P_{\lambda}$.\\
The finite spectral functional calculus is $\mathbb{C}^{\Omega_T}\to\LL(\mathcal{V})$, where $f\mapsto f(T)\defeq\sum\limits_{\lambda\in\Omega_T}f(\lambda)P_{\lambda}$.\\
If $p\in\mathbb{C}[z]$ and $p|_{\Omega_T}=f$, then $p(T)=f(T)$. Hence every $f\in\mathbb{C}^{\Omega_T}$ has a polynomial representative in $T$.
\begin{proof}
For $p|_{\Omega_T}=f$, Corollary \ref{cor:normal-spectral-resolution-functional-calculus} gives $p(T)=\sum\limits_{\lambda\in\Omega_T}p(\lambda)P_{\lambda} =\sum\limits_{\lambda\in\Omega_T}f(\lambda)P_{\lambda} =f(T)$.\\
Since $\Omega_T$ is finite, Lagrange interpolation gives $p\in\mathbb{C}[z]$ with $p|_{\Omega_T}=f$.\\
\hlt{Structural Commentary}: Spectral sums reduce finite functions to polynomial evaluation.\\
\hlt{Key Insight}: Finite spectral data is interpolated by one polynomial.
\end{proof}
\end{theorem}

\begin{theorem} \label{thm:normal-spectral-functional-calculus-rules} \hlt{Rules for Normal Spectral Functions}\\
Let $T\in\LL(\mathcal{V})$ be normal, with spectral resolution $T=\sum\limits_{\lambda\in\Omega_T}\lambda P_{\lambda}$. For $f,g:\Omega_T\to\mathbb{C}$:
\begin{enumerate}[label=(\roman*)]
\setlength{\itemsep}{0pt}
\item $(f+g)(T)=f(T)+g(T)$ and $(fg)(T)=f(T)g(T)$;
\item $f(T)^*=\overline{f}(T)$, where $\overline{f}(\lambda)=\overline{f(\lambda)}$;
\item $f(T)$ is normal and $\operatorname{spec}(f(T))=f(\Omega_T)$.
\end{enumerate}
\begin{proof}
\begin{enumerate}[label=(\roman*)]
\setlength{\itemsep}{0pt}
\item Since $P_{\lambda}P_{\mu}=\delta_{\lambda\mu}P_{\lambda}$, $f(T)g(T)=\sum\limits_{\lambda\in\Omega_T}f(\lambda)g(\lambda)P_{\lambda}=(fg)(T)$; and $(f+g)(T)=f(T)+g(T)$.
\item Since $P_{\lambda}^*=P_{\lambda}$, $f(T)^*=\sum\limits_{\lambda\in\Omega_T}\overline{f(\lambda)}P_{\lambda}=\overline{f}(T)$.
\item For each $\lambda\in\Omega_T$, $f(T)|_{\operatorname{range}(P_{\lambda})}=f(\lambda)I$.\\
Thus $f(T)$ is diagonal in an orthonormal eigenbasis, hence normal, and $\operatorname{spec}(f(T))=f(\Omega_T)$.
\end{enumerate}
\hlt{Structural Commentary}: Spectral projections make operator rules pointwise.\\
\hlt{Key Insight}: Spectral functions act as scalars on spectral subspaces.
\end{proof}
\end{theorem}

\begin{definition} \label{def:positive-operator} \hlt{Positive Operator}\\
$T\in\LL(\mathcal{V})$ is \hlt{positive} if $T=T^*$ and $\innerprod{T\mathbf{v},\mathbf{v}}\geq0$ for all $\mathbf{v}\in\mathcal{V}$.
\end{definition}

\begin{definition} \label{def:square-root-operator} \hlt{Square Root of an Operator}\\
An operator $R\in\LL(\mathcal{V})$ is a \hlt{square root} of $T\in\LL(\mathcal{V})$ if $R^2=T$.
\end{definition}

\begin{theorem} \label{thm:positive-operator-characterization} \hlt{Characterization of Positive Operators}\\
For $T\in\LL(\mathcal{V})$, the following are equivalent:
\begin{enumerate}[label=\roman*.]
\setlength{\itemsep}{0pt}
\item $T$ is positive;
\item $T$ is self-adjoint and every eigenvalue of $T$ is nonnegative;
\item $T$ has a positive square root;
\item $T$ has a self-adjoint square root;
\item there exists $R\in\LL(\mathcal{V})$ such that $T=R^*R$.
\end{enumerate}
\begin{proof}
\begin{enumerate}[label=(\roman*).]
\setlength{\itemsep}{0pt}
\item[$(i)$]$\Rightarrow (ii)$ If $T\mathbf{v}=\lambda\mathbf{v}$ and $\mathbf{v}\neq\mathbf{0}$, then $\lambda\in\mathbb{R}$ by Theorem \ref{thm:self-adjoint-eigenvalues-real}, and $0\leq\innerprod{T\mathbf{v},\mathbf{v}}=\lambda\norm{\mathbf{v}}^2$, so $\lambda\geq0$.
\item[$(ii)$]$\Rightarrow (iii)$ By the real or complex spectral theorem, choose an orthonormal eigenbasis $(\mathbf{e}_1,\ldots,\mathbf{e}_n)$ with $T\mathbf{e}_j=\lambda_j\mathbf{e}_j$ and $\lambda_j\geq0$. Define $S\mathbf{e}_j=\sqrt{\lambda_j}\mathbf{e}_j$ and extend linearly.\\
Then $S$ is diagonal with nonnegative real diagonal entries in an orthonormal basis, hence self-adjoint and positive, and $S^2\mathbf{e}_j=\lambda_j\mathbf{e}_j=T\mathbf{e}_j$ for every $j$. Thus $S^2=T$.
\item[$(iii)$]$\Rightarrow (iv)$ A positive square root is self-adjoint.
\item[$(iv)$]$\Rightarrow (v)$ If $T=R^2$ with $R=R^*$, then $T=R^*R$.
\item[$(v)$]$\Rightarrow (i)$ If $T=R^*R$, then $T^*=T$ and $\innerprod{T\mathbf{v},\mathbf{v}}=\innerprod{R\mathbf{v},R\mathbf{v}}=\norm{R\mathbf{v}}^2\geq0$ for every $\mathbf{v}$, so $T$ is positive.
\end{enumerate}
\hlt{Structural Commentary}: Positivity equals nonnegative spectrum and square-root factorization.\\
\hlt{Key Insight}: In an eigenbasis, the proof is scalar: $\lambda\geq0$ iff $\sqrt{\lambda}$ exists.
\end{proof}
\end{theorem}

\begin{proposition} \label{prop:positive-square-root-unique} \hlt{Unique Positive Square Root}\\
Every positive operator on $\mathcal{V}$ has a unique positive square root.
\begin{proof}
Existence follows from Theorem \ref{thm:positive-operator-characterization}. For uniqueness, decompose $\mathcal{V}=E_{\lambda_1}(T)\oplus\cdots\oplus E_{\lambda_m}(T)$.\\
Let $R$ be a positive square root of $T$. Since $R^2=T$, we have $RT=R^3=TR$, so each $E_{\lambda_j}(T)$ is invariant under $R$. On $E_{\lambda_j}(T)$, the restriction $R_j\defeq R|_{E_{\lambda_j}(T)}$ is positive and satisfies $R_j^2=\lambda_jI$.\\
By the spectral theorem for the positive operator $R_j$, every eigenvalue $\rho$ of $R_j$ satisfies $\rho\geq0$ and $\rho^2=\lambda_j$, hence $\rho=\sqrt{\lambda_j}$. Therefore $R_j=\sqrt{\lambda_j}I$ on $E_{\lambda_j}(T)$.\\
Thus every positive square root agrees with the same scalar operator on each eigenspace $E_{\lambda_j}(T)$, so it is unique.\\
\hlt{Structural Commentary}: Each eigenspace determines its own scalar root.\\
\hlt{Key Insight}: Positivity selects the nonnegative root.
\end{proof}
\end{proposition}

\begin{definition} \label{def:positive-square-root-notation} \hlt{Positive Square Root Notation}\\
If $T$ is positive, then $\sqrt{T}$ denotes the unique positive square root of $T$.
\end{definition}

\begin{example} \label{ex:spectral-functional-calculus-basic-functions} \hlt{Basic Functions from the Spectral Calculus}\\
Let $T\in\LL(\mathcal{V})$ be normal, with spectral resolution $T=\sum\limits_{\lambda\in\Omega_T}\lambda P_{\lambda}$. Then $T^*=\sum\limits_{\lambda\in\Omega_T}\overline{\lambda}P_{\lambda}$.\\
If $0\notin\Omega_T$, then $T^{-1}=\sum\limits_{\lambda\in\Omega_T}\lambda^{-1}P_{\lambda}$. If $T\geq0$, then $T^{1/2}=\sum\limits_{\lambda\in\Omega_T}\lambda^{1/2}P_{\lambda}$, where $\lambda^{1/2}\geq0$.
\end{example}

\begin{theorem} \label{thm:normal-spectral-projection-commutativity} \hlt{Commutativity via Spectral Projections}\\
Let $T,S\in\LL(\mathcal{V})$ be normal, with spectral resolutions $T=\sum\limits_{\lambda\in\Omega_T}\lambda P_{\lambda}$ and $S=\sum\limits_{\mu\in\Omega_S}\mu Q_{\mu}$.
\begin{enumerate}[label=(\roman*)]
\setlength{\itemsep}{0pt}
\item For $R\in\LL(\mathcal{V})$, $RT=TR$ if and only if $RP_{\lambda}=P_{\lambda}R$ for all $\lambda\in\Omega_T$.
\item $TS=ST$ if and only if $P_{\lambda}Q_{\mu}=Q_{\mu}P_{\lambda}$ for all $\lambda\in\Omega_T$ and $\mu\in\Omega_S$.
\item If $f:\Omega_T\to\mathbb{C}$ and $g:\Omega_S\to\mathbb{C}$ are injective, then $f(T)g(S)=g(S)f(T)$ if and only if $TS=ST$.
\end{enumerate}
\begin{proof}
\begin{enumerate}[label=(\roman*).]
\setlength{\itemsep}{0pt}
\item If $RT=TR$, then $Rp(T)=p(T)R$ for all $p\in\mathbb{C}[z]$. Since $P_{\lambda}\in\mathbb{C}[T]$ by Corollary \ref{cor:normal-spectral-resolution-functional-calculus}, $RP_{\lambda}=P_{\lambda}R$. Conversely, if $RP_{\lambda}=P_{\lambda}R$ for all $\lambda\in\Omega_T$, then $RT=\sum\limits_{\lambda\in\Omega_T}\lambda RP_{\lambda} =\sum\limits_{\lambda\in\Omega_T}\lambda P_{\lambda}R =TR$.
\item If $TS=ST$, then (i) with $R=S$ gives $SP_{\lambda}=P_{\lambda}S$. Since $Q_{\mu}\in\mathbb{C}[S]$, $P_{\lambda}Q_{\mu}=Q_{\mu}P_{\lambda}$ for all $\lambda,\mu$. Conversely, if $P_{\lambda}Q_{\mu}=Q_{\mu}P_{\lambda}$ for all $\lambda,\mu$, then $TS=\sum\limits_{\lambda,\mu}\lambda\mu P_{\lambda}Q_{\mu} =\sum\limits_{\lambda,\mu}\lambda\mu Q_{\mu}P_{\lambda} =ST$
\item Suppose $f$ and $g$ are injective. If $f(T)g(S)=g(S)f(T)$, choose $h:f(\Omega_T)\to\Omega_T$ and $k:g(\Omega_S)\to\Omega_S$ with $h\circ f=\operatorname{id}_{\Omega_T}$ and $k\circ g=\operatorname{id}_{\Omega_S}$. By Theorem \ref{thm:finite-spectral-functional-calculus-normal}, $T=h(f(T))$ and $S=k(g(S))$, so $TS=ST$. Conversely, if $TS=ST$, then (ii) gives $P_{\lambda}Q_{\mu}=Q_{\mu}P_{\lambda}$,\\
and $f(T)g(S)=\sum\limits_{\lambda,\mu}f(\lambda)g(\mu)P_{\lambda}Q_{\mu} =\sum\limits_{\lambda,\mu}f(\lambda)g(\mu)Q_{\mu}P_{\lambda} =g(S)f(T)$.
\end{enumerate}
\hlt{Structural Commentary}: Spectral projections are the coordinate idempotents of a normal operator.\\
\hlt{Key Insight}: To commute with a normal operator is to preserve each of its eigenspaces.
\end{proof}
\end{theorem}

\begin{corollary} \label{cor:commuting-normal-pair-generated-by-one-normal-operator} \hlt{A Commuting Normal Pair Comes from One Normal Operator}\\
Let $T,S\in\LL(\mathcal{V})$ be normal on a finite-dimensional complex inner-product space.\\
Then $TS=ST$ $\Longleftrightarrow$ $\exists\, H\in\LL(\mathcal{V})$ normal, $\ p,q\in\mathbb{C}[z]$ such that $T=p(H),\ S=q(H)$.\\
If $TS=ST$, then $H$ may be chosen in $\mathbb{C}[T,S]$.
\begin{proof}
\begin{enumerate}[label=\roman*.]
\setlength{\itemsep}{0pt}
\item[$(\Rightarrow)$] Suppose $TS=ST$. Write $T=\sum\limits_{\lambda\in\Omega_T}\lambda P_{\lambda}$ and $S=\sum\limits_{\mu\in\Omega_S}\mu Q_{\mu}$.\\
By Theorem \ref{thm:normal-spectral-projection-commutativity}, $P_{\lambda}Q_{\mu}=Q_{\mu}P_{\lambda}$. Let $\mathcal{I}=\{(\lambda,\mu)\in\Omega_T\times\Omega_S:P_{\lambda}Q_{\mu}\neq0\}$.\\
The projections $P_{\lambda}Q_{\mu}$, $(\lambda,\mu)\in\mathcal{I}$, are pairwise orthogonal and sum to $I_{\mathcal{V}}$.\\
Choose distinct scalars $\alpha_{\lambda\mu}$ for $(\lambda,\mu)\in\mathcal{I}$ and set $H=\sum\limits_{(\lambda,\mu)\in\mathcal{I}}\alpha_{\lambda\mu}P_{\lambda}Q_{\mu}$.\\
Then $H$ is normal, and $H\in\mathbb{C}[T,S]$ since $P_{\lambda}\in\mathbb{C}[T]$ and $Q_{\mu}\in\mathbb{C}[S]$.\\
Choose $p,q\in\mathbb{C}[z]$ with $p(\alpha_{\lambda\mu})=\lambda$ and $q(\alpha_{\lambda\mu})=\mu$ on $\mathcal{I}$.\\
Hence $p(H)=\sum\limits_{(\lambda,\mu)\in\mathcal{I}}\lambda P_{\lambda}Q_{\mu}=T$, and $q(H)=\sum\limits_{(\lambda,\mu)\in\mathcal{I}}\mu P_{\lambda}Q_{\mu}=S$.
\item[$(\Leftarrow)$] If $T=p(H)$ and $S=q(H)$, then $TS=p(H)q(H)=q(H)p(H)=ST$.
\end{enumerate}
\hlt{Structural Commentary}: One scalar labels each common spectral block.\\
\hlt{Key Insight}: Joint spectral projections reduce a pair of commuting normal operators to finite scalar data.
\end{proof}
\end{corollary}

\begin{definition} \label{def:matrix-exponential} \hlt{Matrix Exponential}\\
For $A\in\operatorname{Mat}(n,n,\mathbb{C})$ and $t\in\mathbb{C}$, define $e^{tA}\defeq\sum\limits_{k=0}^{\infty}\frac{t^kA^k}{k!}$.
\end{definition}

\begin{proposition} \label{prop:diagonalizable-pure-modes-powers-exponentials} \hlt{Pure Modes for Powers and Exponentials}\\
Let $A\in\operatorname{Mat}(n,n,\mathbb{C})$ be diagonalizable with $A=S\Lambda S^{-1}$, where $\Lambda=\operatorname{diag}(\lambda_1,\ldots,\lambda_n)$ and the columns of $S$ are $\mathbf{v}_1,\ldots,\mathbf{v}_n$. If $\mathbf{x}_0=\sum\limits_{j=1}^n c_j\mathbf{v}_j$, then $A^k\mathbf{x}_0=\sum\limits_{j=1}^n c_j\lambda_j^k\mathbf{v}_j$ for $k\geq0$; and $e^{tA}\mathbf{x}_0=\sum\limits_{j=1}^n c_j e^{\lambda_jt}\mathbf{v}_j$ for $t\in\mathbb{C}$.
\begin{proof}
Let $c=(c_1,\ldots,c_n)^t$, so $\mathbf{x}_0=Sc$.\\
Then $A^k=S\Lambda^kS^{-1}$, and Definition \ref{def:matrix-exponential} gives $e^{tA}=S\left(\sum\limits_{r=0}^{\infty}\frac{t^r\Lambda^r}{r!}\right)S^{-1} =S\operatorname{diag}(e^{\lambda_1t},\ldots,e^{\lambda_nt})S^{-1}$.\\
Applying these two formulas to $Sc$ gives the stated identities.\\
\hlt{Structural Commentary}: Eigenbasis coordinates decouple powers and exponentials.\\
\hlt{Key Insight}: Diagonalization turns coupled evolution into independent scalar modes.
\end{proof}
\end{proposition}
